% SPDX-License-Identifier: Apache-2.0
\section{Where the quantum method sits, and why there}
\label{sec:hybrid}

\textbf{Classical, on all traffic.} A gradient-boosted scorer runs on every transaction --- the
only component that sees full volume, and deliberately the strongest available baseline at
\ClaimGbdtFullAucMin--\ClaimGbdtFullAucMax{} ROC AUC over five seeds. A hybrid whose classical
half is weak proves nothing about its quantum half.

\textbf{Conformal, on the decision.} The certificate wraps whatever scores the band, so the
deliverable an issuer can act on survives a negative quantum result --- as it had to here.

\textbf{Quantum, on the band only.} At a \ClaimTestFraudRate\,\% positive rate and a 3-D Secure
challenge budget the band holds thousands of rows rather than hundreds of thousands, and a
kernel whose Gram matrix costs $O(n^2)$ circuit evaluations is affordable on thousands. The band
is not a convenience; it is the only place in this pipeline where the arithmetic works.

\textbf{The compute path.} Both arms run on simulators, which the statement permits: no number
here comes from a QPU. The fidelity kernel's reference implementation is an exact statevector,
with Braket's \texttt{LocalSimulator} and Aer on CPU and GPU agreeing to within
\ClaimParityWorstDifference{} over \ClaimParityComparisons{} comparisons; the tensor network
runs in PyTorch on one GPU. On the sample count the statement asks for: the screens ran on
\ClaimScreenSamples{} stratified rows per configuration across \ClaimScreenConfigs{}
configurations, and the kernel arm was rejected before any downstream evaluation, so no larger
quantum execution took place. The circuits are within near-term reach --- at most
\ClaimCircuitQubitsMax{} qubits, transpiled depth \ClaimCircuitDepthMax{} and
\ClaimCircuitTwoQubitMax{} two-qubit gates on a portable basis --- so what stopped this arm was
the screens, not the hardware.
